\documentclass[11pt,a4paper]{article}

% ==================================================
% VARIABLES – adjust per lab
% ==================================================
\newcommand{\course}{Software and Process Notations}
\newcommand{\labno}{1}
\newcommand{\labtopic}{System Sketches}
\newcommand{\reportdate}{21 October 2026}
\newcommand{\headernames}{Stirbu, Hosseini, Stolzlechner}
\newcommand{\pdfauthorname}{Ricardo Stolzlechner}

% Laborantenliste: \laborant{Name}{E-Mail} (eine Zeile pro Person)
\newcommand{\members}{%
    \member{Alex Stirbu}{astirbu-blagut.cysb-m2026@fh-salzburg.ac.at}    
    \member{Manizheh Hosseini}{mhosseini.cysb-m2026@fh-salzburg.ac.at}
    \member{Ricardo Stolzlechner}{ricardo@stolzlechner.at}
}

% Verwendetes KI-Werkzeug (laut KI-Richtlinie 1.0, Kap. 1.7)
\newcommand{\aitool}{Claude Code (Anthropic, Claude Opus 5.5, token based billing)}

% --------------------------------------------------
% Packages
% --------------------------------------------------
\usepackage[utf8]{inputenc}
\usepackage[T1]{fontenc}
\usepackage[english]{babel}
\usepackage[a4paper,margin=2.5cm,headheight=14pt]{geometry}

\usepackage{amsmath, amssymb}
\usepackage{enumitem}
\usepackage{xcolor}
\usepackage{fancyhdr}
\usepackage{array}
\usepackage{booktabs}
\usepackage{graphicx}
\usepackage{listings}
\usepackage{setspace}
\usepackage{hyperref}

% Listings
\definecolor{codegreen}{rgb}{0,0.6,0}
\definecolor{codegray}{rgb}{0.5,0.5,0.5}
\definecolor{codepurple}{rgb}{0.58,0,0.82}
\definecolor{backcolour}{rgb}{0.95,0.95,0.92}

\lstdefinestyle{mystyle}{
    backgroundcolor=\color{backcolour},
    commentstyle=\color{codegreen},
    keywordstyle=\color{magenta},
    numberstyle=\tiny\color{codegray},
    stringstyle=\color{codepurple},
    basicstyle=\ttfamily\footnotesize,
    breaklines=true,
    captionpos=b,
    keepspaces=true,
    numbers=left,
    numbersep=5pt,
    showstringspaces=false,
    tabsize=2
}
\lstset{style=mystyle}

% Title page
\newcommand{\member}[2]{#1 & #2 \\}

% Spacing
\onehalfspacing
\setlength{\parindent}{0pt}
\setlength{\parskip}{6pt}

% --------------------------------------------------
% Header and footer
% --------------------------------------------------
\pagestyle{fancy}
\fancyhf{}
\lhead{\course}
\rhead{\headernames}
\cfoot{Page \thepage\ of \pageref{LastPage}}

% --------------------------------------------------
% Hyperref settings
% --------------------------------------------------
\hypersetup{
    colorlinks=true,
    linkcolor=blue,
    urlcolor=blue,
    pdftitle={\course\ -- Lab \labno: \labtopic},
    pdfauthor={\pdfauthorname}
}

% ==================================================
% Document
% ==================================================
\begin{document}

% --------------------------------------------------
% Page 1: Title page
% --------------------------------------------------
\pagenumbering{Alph}
\begin{titlepage}
    \centering
    \vspace*{4cm}
    {\LARGE \textbf{Lab Report - Lab Unit \labno}}\\[0.5em]
    {\Large \labtopic}\\[0.5em]
    {\large \course}

    \vspace{3cm}

    \begin{tabular}{l l}
        \toprule
        \textbf{Lab Member} & \textbf{E-mail} \\
        \midrule
        \members
        \bottomrule
    \end{tabular}

    \vfill
    {\large \reportdate}
\end{titlepage}
\pagenumbering{arabic}

% --------------------------------------------------
% Page 2: Content
% --------------------------------------------------
\section*{Objective/Physical-Based}

\begin{figure}[htbp]
    \centering
    \includegraphics[width=\linewidth]{img/physical.png}
    \caption{Objective/physical-based view: components of the system and their interactions}
    \label{fig:physical}
\end{figure}

\newpage

\section*{Process-Based}

\begin{figure}[htbp]
    \centering
    \includegraphics[width=\linewidth,height=0.8\textheight,keepaspectratio]{img/functional.png}
    \caption{Process-based view: sequence of actions incl.\ alternate courses and exceptions}
    \label{fig:functional}
\end{figure}

\newpage

\section*{Time-Based}

\begin{figure}[htbp]
    \centering
    \includegraphics[width=\linewidth]{img/time.png}
    \caption{Time-based view: system behaviour and data lifetime over one parking session}
    \label{fig:time}
\end{figure}


% --------------------------------------------------
% Page 3: AI disclosure
% --------------------------------------------------
\clearpage
\section*{Statement on the Use of AI}
\addcontentsline{toc}{section}{Statement on the Use of AI}

In preparing this lab report, \aitool{} was used as a generative AI tool
exclusively to assist with the following tasks:

\begin{itemize}
    \item creating and designing the \LaTeX{} layout,
    \item drafting the three system sketches (draw.io),
    \item checking spelling and grammar,
    \item research.
\end{itemize}

The \LaTeX{} template of this report, including this statement, was created
with the assistance of the AI tool and reviewed and adapted by the lab members.
First drafts of the three system sketches were generated by the AI tool based on
the use case description written by the lab members; the lab members subsequently
reviewed, revised and finalized the sketches. Apart from that, all content,
results and conclusions were developed independently; no AI-generated text was
adopted. The tool was used in line with
the principles of \emph{responsible AI use} (transparency, verifiability, human
oversight). Sources found through research were checked for accuracy
independently. All content was reviewed, validated and finally approved by the
lab members (cf.\ FH Salzburg AI guideline \emph{KI-Richtlinie}, version~1.0,
ch.~1.3 and 2.1.2).

\label{LastPage}
\end{document}
